\section{Smartphone activity recognition on held-out subjects (UCI HAR)}
\label{app:har}

\subsection{Split, representation summary, and fitting protocol}
The UCI Human Activity Recognition (HAR) benchmark~\cite{har} records six
smartphone activities (walking, walking upstairs, walking downstairs, sitting,
standing, and laying) at \(50\,\mathrm{Hz}\) in fixed windows of \(R = 128\)
samples (\(2.56\,\mathrm{s}\)). We use the official subject-disjoint partition:
\(7{,}352\) training windows from \(21\) participants and \(2{,}947\) test
windows from \(9\) held-out participants. Each window supplies three-axis body
acceleration \(b_r \in \R^3\), angular velocity \(\omega_r \in \R^3\), and total
acceleration \(a_r \in \R^3\), from which we form the gravitational component
\(g_r = a_r - b_r\). No additional filtering or re-segmentation is applied.

Recognizing both static postures and dynamic gaits with a single linear readout
calls for complementary views of each window: gravitational orientation and
ordered cross-axis motion in the time domain, oscillation cadence and harmonic
spread in the frequency domain, and frame-invariant magnitude, alignment, and
short-lag persistence. The \emph{compact368} representation concatenates three
corresponding branches whose feature counts are summarized up front in
Table~\ref{tab:compact-blocks}.

\begin{table}[ht]
\centering
\begin{tabular}{@{}lrr@{}}
\toprule
\emph{compact368} component & Nominal & Retained \\
\midrule
Raw sensor \(P=3\) QR-Pascal depth-two logsignature (\(18\) channels) & 171 & 171 \\
Spectral shape (\(20\) scalar \(P=3\) QR-Pascal depth-two logsignatures) & 120 & 100 \\
Spectral scale (mean and std of unnormalized bin amplitudes, \(20\) signals) & 40 & 40 \\
Quadratic geometric streams (mean, std, skewness of \(12\) streams) & 36 & 36 \\
Instantaneous cosine alignment (mean and std of \(3\) cross-triad streams) & 6 & 6 \\
Window-level linear trends (\(q_1\)-weighted trend across all \(15\) streams) & 15 & 15 \\
\midrule
Total & 388 & 368 \\
\bottomrule
\end{tabular}
\caption{Feature breakdown of \emph{compact368} before and after applying the
training-only variance mask, which removes the \(20\) constant unit-mass
spectral endpoints.}
\label{tab:compact-blocks}
\end{table}

Every experiment on the official UCI HAR split in this section uses an identical,
strictly training-fitted preprocessing and classification pipeline:
\begin{enumerate}
\item \emph{Sensor calibration:} All per-axis training means and standard
deviations (raw and spectral branches) and per-triad root-mean-square scales
(geometric branch) are estimated across the \(7{,}352\) training windows only
and applied unchanged to test windows.
\item \emph{Constant-column variance mask:} Before fitting the classifier,
feature column \(j\) is discarded if its training population standard deviation
satisfies \(\sigma^{\mathrm{train}}_j \leq 10^{-6}\max_i |F^{\mathrm{train}}_{ij}|\).
On HAR this removes precisely the \(20\) unit-sum spectral endpoints, leaving
\(368\) columns.
\item \emph{Feature standardization and ridge readout:} Each retained column is
standardized using its training mean and training population standard deviation.
Writing \(F \in \R^{N_{\mathrm{train}} \times d}\) for the standardized feature
matrix and \(Y \in \{\pm 1\}^{N_{\mathrm{train}} \times 6}\) for one-versus-rest
target indicators (\(+1\) for the target activity, \(-1\) otherwise), we solve
the closed-form ridge problem with an unpenalized intercept \(b \in \R^6\) and
fixed penalty \(\alpha = 10\):
\begin{equation}\label{eq:har-ridge}
 \min_{b \in \R^6,\, W \in \R^{d \times 6}}
 \|Y - \mathbf 1 b^\top - F W\|_F^2 + \alpha \|W\|_F^2.
\end{equation}
Test windows are standardized with the frozen training parameters and assigned
to \(\arg\max_{c} \bigl(b_c + (F_{\mathrm{test}} W)_{ic}\bigr)\).
\end{enumerate}

\subsection{Specification of the three feature branches}

\paragraph{1. Raw sensor branch: posture levels and temporal order (171 features).}
To preserve static gravitational orientation across windows, we take the six raw
channels \((a_r, \omega_r) \in \R^6\) and standardize each axis using the global
mean and population standard deviation computed across all \(7{,}352 \times 128\)
training samples, \emph{without} centering individual windows. Denoting the
standardized reading at sample \(r \in \{0, \ldots, 127\}\) by \(x_r \in \R^6\),
we treat scaled readings \(x_r / 128\) as increments and lift them with the
\(P = 3\) orthonormal QR-Pascal weights \(q_p(r)\) of Section~\ref{app:qr}:
\[
 v_r^{(c, p)} = \frac{x_r^c\,q_p(r)}{128},
 \qquad c \in \{1, \ldots, 6\},\quad p \in \{0, 1, 2\}.
\]
Lifted coordinates are ordered by sensor axis \(c\) first and QR copy \(p\)
second, and accumulated from \(y_0 = 0\) without an added clock channel. Because
\(q_0(r) = 1/\sqrt{128}\), the constant-copy endpoint equals
\(\sum_{r=0}^{127} x_r^c / 128^{3/2} = \bar x^c / \sqrt{128}\), retaining the
window mean alongside linear and quadratic temporal moments and ordered
cross-channel areas. The depth-two logsignature of this \(18\)-channel path
contributes \(18\) endpoints and \(\binom{18}{2} = 153\) signed areas
\(J_{ij} = \tfrac12(B_{ij} - B_{ji})\), totaling \(171\) features.

\paragraph{2. Spectral branch: frequency shape and amplitude scale (140 retained features).}
To summarize rhythmic cadence and harmonic spread, we start from five
three-axis triads: \(b, g, \omega\) (length \(n = 128\)) and the
unscaled first differences \(\Delta b_r = b_{r+1} - b_r\) and
\(\Delta \omega_r = \omega_{r+1} - \omega_r\) (length \(n = 127\)). Each triad is
standardized per axis using training-set statistics as in the raw branch, after
which we append its instantaneous Euclidean norm \(\|x_r\|\) \emph{in the time
domain} to obtain \(20\) scalar signals (\(12\) of length \(128\) and \(8\) of
length \(127\)). Forming \(\|x_r\|\) prior to the Fourier transform captures
nonlinear envelope modulations; its power spectrum is distinct from the sum of
the three individual axis spectra. The power spectrum of each resulting scalar
signal is invariant under circular shifts of that signal. A circular shift of
the raw window can change its first-difference spectra, since differencing
omits the wraparound step.

For each scalar signal \(x_0, \ldots, x_{n-1}\), we compute the Parseval-scaled
one-sided discrete Fourier power spectrum over \(K = \lfloor n/2 \rfloor + 1\)
bins (\(K = 65\) for \(n = 128\), \(K = 64\) for \(n = 127\)):
\[
 \widehat x_k = \sum_{r=0}^{n-1} x_r e^{-2\pi i k r / n}, \qquad
 e_k = w_k \left|\frac{\widehat x_k}{n}\right|^2, \qquad
 E = \sum_{k=0}^{K-1} e_k = \frac{1}{n}\sum_{r=0}^{n-1} x_r^2,
\]
where \(w_k = 1\) at DC (\(k = 0\)) and, when \(n\) is even, at the Nyquist bin
(\(k = n/2\)), and \(w_k = 2\) at all other bins, including the highest bin for
odd \(n\). Normalized power shares are \(p_k = e_k / E\), and unnormalized bin
amplitudes are
\(a_k = \sqrt{e_k}\). Treating each signal's shares \(p_k\) independently as
scalar increments ordered by bin \(k = 0, \ldots, K-1\) and lifting with the
three QR weights \((q_0(k), q_1(k), q_2(k))\) gives \(3 + \binom{3}{2} = 6\)
depth-two logsignature coordinates per spectrum, with the three areas stored at
scale \(\sqrt 2\,J_{ij} = (B_{ij} - B_{ji})/\sqrt 2\) as analyzed
in~\eqref{eq:spectral-qr-scale}. Because \(\sum_{k=0}^{K-1} p_k = 1\), the first
endpoint \(q_0(0)\sum_k p_k = 1/\sqrt K\) is constant across windows and removed
by the variance mask, leaving \(20 \times 5 = 100\) nonconstant spectral-shape
features (with no cross-spectrum areas). Appending the mean and population
standard deviation of \(a_0, \ldots, a_{K-1}\) for each signal adds
\(20 \times 2 = 40\) amplitude-scale features.

\paragraph{3. Geometric branch: magnitude, persistence, and alignment (57 features).}
To capture sensor intensity, inter-triad alignment, and short-lag autocorrelation
independently of coordinate-axis sign conventions, we calibrate the uncentered
physical triads \(b, g, \omega\) by dividing each triad by its scalar root mean
square over all three axes, time samples, and training windows. Using one scalar
RMS per triad rather than per-axis centering ensures that the calibrated vectors
\(u_1(r), u_2(r), u_3(r) \in \R^3\) transform covariantly under any constant
orthogonal rotation \(O \in \mathrm{O}(3)\) applied simultaneously to all three
triads. From \(u_1, u_2, u_3\) we form \(12\) rotation-invariant quadratic
streams,
\[
 \|u_i(r)\|^2 \quad (1 \leq i \leq 3), \qquad
 \langle u_i(r), u_j(r)\rangle \quad (1 \leq i < j \leq 3), \qquad
 \langle u_i(r), u_i(r - \tau)\rangle \quad (1 \leq i \leq 3;\ \tau \in \{2, 5\}),
\]
where lags \(\tau = 2\) and \(\tau = 5\) correspond to \(40\,\mathrm{ms}\) and
\(100\,\mathrm{ms}\) (giving valid lengths \(L = 128, 126, 123\)), together with
three instantaneous cosine-similarity streams
\(\langle u_i(r), u_j(r)\rangle / (\|u_i(r)\|\|u_j(r)\|)\) (\(1 \leq i < j \leq 3\)).
Within each window, every stream \(z_0, \ldots, z_{L-1}\) contributes its mean
\(\bar z\) and population standard deviation \(\sigma_z\); each of the \(12\)
quadratic streams additionally contributes its skewness
\(L^{-1}\sum_{r=0}^{L-1} y_r^3\), where \(y_r = (z_r - \bar z)/\sigma_z\); and all
\(15\) streams contribute their linear QR-Pascal trend
\[
 T(z) = \frac{1}{L}\sum_{r=0}^{L-1} q_1(r) y_r,
\]
which is positive when standardized stream values increase toward the end of the
window. Taken alone, these \(36 + 6 + 15 = 57\) features constitute the
\emph{geometry57} representation.

\begin{remark}[Numerical safeguards and archived floating-point conventions]
To guard against degenerate denominators without altering well-conditioned
windows, per-axis training standard deviations and per-triad training RMS scales
are floored at \(10^{-8}\); spectral total power uses \(\max(E, 10^{-20})\),
with a unit-mass check after normalization that rejects zero-power inputs
(all evaluated HAR spectra summed to one to machine precision);
cosine-similarity denominators use \(\max(\|u_i\|\|u_j\|, 10^{-12})\);
and within-window standardized streams set \(y_r = 0\) identically whenever
\(\sigma_z^2 \leq 10^{-12}\) (otherwise flooring \(\sigma_z\) at \(10^{-8}\)). To
follow the archived rounding convention, raw and spectral signature tensors
are accumulated in binary64 and cast to binary32 immediately before extracting
upper-triangular areas, whereas geometric summaries and the ridge solve remain in
binary64 (\nolinkurl{reproduction/uci_har.py}).
\end{remark}

\subsection{Standalone blocks, supplied baseline, and single-block deletions}
\label{app:ablations}
On the \(2{,}947\) held-out subject test windows, \emph{compact368} achieves
\(98.03\%\) accuracy (\(58\) errors). Two complementary ablations clarify how its
components contribute to the \(\alpha = 10\) linear readout: fitting the readout
to individual blocks in isolation (Table~\ref{tab:har-standalone}), and removing
one block at a time from the full representation (Table~\ref{tab:har-deletions}).

\begin{table}[ht]
\centering
\begin{tabular}{@{}lrrr@{}}
\toprule
Standalone representation or baseline & Features & Accuracy (\%) & Errors \\
\midrule
Raw QR-Pascal logsignature only & 171 & 81.98 & 531 \\
Spectral QR-Pascal shape only (no amplitude scale) & 100 & 82.90 & 504 \\
Geometric branch (\emph{geometry57}) only & 57 & 85.14 & 438 \\
Supplied UCI HAR hand-engineered baseline & 561 & 95.72 & 126 \\
Full \emph{compact368} & 368 & 98.03 & 58 \\
\bottomrule
\end{tabular}
\caption{Standalone branch evaluations and the 561-feature benchmark baseline on
the official subject-disjoint UCI HAR split (\(2{,}947\) test windows) under the
identical \(\alpha = 10\) ridge readout.}
\label{tab:har-standalone}
\end{table}

\begin{table}[ht]
\centering
\begin{tabular}{@{}lrrrr@{}}
\toprule
Leave-one-out deletion from \emph{compact368} & Retained features & Accuracy (\%) & Errors & \(\Delta\) Errors \\
\midrule
All blocks retained (\emph{compact368}) & 368 & 98.03 & 58 & \(0\) \\
Remove raw QR-Pascal branch (\(171\)) & 197 & 96.17 & 113 & \(+55\) \\
Remove spectral QR-Pascal shape (\(100\)) & 268 & 96.78 & 95 & \(+37\) \\
Remove spectral amplitude mean/std (\(40\)) & 328 & 97.15 & 84 & \(+26\) \\
Remove cross-triad cosine mean/std (\(6\)) & 362 & 96.34 & 108 & \(+50\) \\
Remove quadratic stream mean/std/skew (\(36\)) & 332 & 97.42 & 76 & \(+18\) \\
Remove geometric/cosine linear trends (\(15\)) & 353 & 97.86 & 63 & \(+5\) \\
\bottomrule
\end{tabular}
\caption{Single-block deletions from \emph{compact368} on the official
subject-disjoint UCI HAR split (\(2{,}947\) test windows).}
\label{tab:har-deletions}
\end{table}

Taken alone, each individual branch makes between \(438\) and \(531\) errors
(\(81.98\%\)--\(85.14\%\)), well below both the 561-feature engineered baseline
(\(95.72\%\), \(126\) errors) and their concatenation (\(98.03\%\), \(58\)
errors). Conversely, Table~\ref{tab:har-deletions} shows that deleting any major
component substantially increases test errors even when all remaining blocks are
present: dropping the raw Pascal branch adds \(55\) errors, dropping the six
cross-triad cosine features adds \(50\) errors, dropping the spectral Pascal
shape coordinates adds \(37\) errors, and dropping the \(40\) unnormalized
amplitude summaries adds \(26\) errors (confirming empirically why unit-normalized
spectra require separate energy summaries). Because ridge weights are refitted
for each subset and correlated signal properties are shared across branches,
these error increments measure conditional usefulness to the linear readout
rather than isolated causal effects. Aggregate error counts are reported because
per-window subject-indexed predictions were not saved in the original
\emph{compact368} archive.

\subsection{Conventional signature replacements and staircase-projection verification}
To test what Pascal increment coding contributes beyond conventional path
signatures within the same pipeline, an archived \(187\)-configuration diagnostic
study replaced either the raw branch or the spectral branch with uncoded
clock-augmented signatures while holding the remaining \emph{compact368} features
and the \(\alpha = 10\) readout fixed.

\paragraph{Replacing the raw sensor branch.}
Keeping the \(140\) spectral and \(57\) geometric features fixed (\(197\)
columns), we replaced the \(171\)-feature raw QR-Pascal block with ordinary
tensor and logsignatures (up to \(1{,}600\) retained raw coordinates) built from
the same six standardized sensor axes. The strongest conventional replacement in
this setting used a clock-augmented, origin-anchored \emph{lead--lag} path: for
each consecutive pair of standardized readings \(x_r, x_{r+1} \in \R^6\), the
\(12\) spatial coordinates visit \((x_r, x_r)\), \((x_{r+1}, x_r)\), and
\((x_{r+1}, x_{r+1})\) in order, augmented by a uniform clock from \(0\) to \(1\)
and an initial step from the origin \(0 \in \R^{13}\). Removing the deterministic
clock endpoint and clock--clock square leaves \(13 + 13^2 - 2 = 180\) depth-two
tensor coordinates. Table~\ref{tab:ordinary-raw} compares this best observed
lead--lag control against \emph{compact368}, alongside a parallel comparison
appended to the \(561\) supplied HAR features (where the conventional control
visits each six-axis reading once with a clock and origin anchor, giving
\(7 + 7^2 - 2 = 54\) coordinates, against all \(18 + 18^2 = 342\) depth-two
tensor coordinates of the \(P = 3\) QR-Pascal lift).

\begin{table}[ht]
\centering
\begin{tabular}{@{}lrrrr@{}}
\toprule
Raw sensor branch representation & Raw & Total & Accuracy (\%) & Errors \\
\midrule
\multicolumn{5}{@{}l}{\emph{Within \emph{compact368} (fixed 140 spectral + 57 geometric features)}} \\
QR-Pascal \(P=3\), depth-two logsignature (\emph{compact368}) & 171 & 368 & 98.03 & 58 \\
Clocked anchored lead--lag tensor signature, depth 2 & 180 & 377 & 97.46 & 75 \\
\midrule
\multicolumn{5}{@{}l}{\emph{Appended to the 561 supplied UCI HAR features}} \\
QR-Pascal \(P=3\), depth-two tensor signature & 342 & 903 & 97.12 & 85 \\
Clocked anchored ordinary tensor signature, depth 2 & 54 & 615 & 97.05 & 87 \\
\bottomrule
\end{tabular}
\caption{Raw-branch replacements on the official subject-disjoint UCI HAR split
under fixed \(\alpha = 10\) ridge classification (\cite[Table~1]{mainpaper}).}
\label{tab:ordinary-raw}
\end{table}

\paragraph{Replacing the spectral branch and checking Theorem~\ref{prop:staircase-projection}.}
Holding the \(171\) raw features, \(57\) geometric features, and \(40\) spectral
amplitude summaries fixed (\(268\) columns), we replaced the \(100\)
QR-Pascal spectral-shape features with ordinary signatures of the two-channel
frequency--mass staircase \((F, M)\) defined in Subsection~\ref{app:staircase-projection},
evaluated independently on each of the \(20\) normalized spectra across depths
\(2\) through \(5\). Because every spectrum has fixed staircase endpoint
\((1, 1)\), words consisting purely of \(F\) letters or purely of \(M\) letters
are constant and discarded (removing both endpoints from the logsignatures).
Table~\ref{tab:ordinary-spectral} reports the resulting dimensions and test
scores, concluding with the explicit linear projection of the depth-five
staircase tensor onto the \(P=3\) QR-Pascal coordinates.

\begin{table}[ht]
\centering
\begin{tabular}{@{}lrrrr@{}}
\toprule
Spectral shape branch & Spectral & Total & Accuracy (\%) & Errors \\
\midrule
QR-Pascal \(P=3\), depth-two logsignature (\emph{compact368}) & 100 & 368 & 98.03 & 58 \\
Ordinary \((F,M)\) staircase logsignature, depth 2 & 20 & 288 & 97.49 & 74 \\
Ordinary \((F,M)\) staircase logsignature, depth 3 & 60 & 328 & 97.79 & 65 \\
Ordinary \((F,M)\) staircase tensor signature, depth 3 & 160 & 428 & 97.83 & 64 \\
Ordinary \((F,M)\) staircase logsignature, depth 4 & 120 & 388 & 97.83 & 64 \\
Ordinary \((F,M)\) staircase tensor signature, depth 4 & 440 & 708 & 98.00 & 59 \\
Ordinary \((F,M)\) staircase logsignature, depth 5 & 240 & 508 & 97.93 & 61 \\
Ordinary \((F,M)\) staircase tensor signature, depth 5 & 1,040 & 1,308 & 97.83 & 64 \\
\midrule
Depth-five staircase projected to \(P=3\) QR-Pascal & 100 & 368 & 98.03 & 58 \\
\bottomrule
\end{tabular}
\caption{Spectral-shape replacements within \emph{compact368} (with the other
\(268\) raw, amplitude-scale, and geometric features held fixed) on the official
UCI HAR split. Counts are retained columns, including linear dependencies among
tensor words. (The depth-two staircase tensor with \(40\) spectral columns also
achieves \(97.49\%\), \(74\) errors.)}
\label{tab:ordinary-spectral}
\end{table}

Table~\ref{tab:ordinary-spectral} illustrates three distinct theoretical points
from Section~\ref{app:depth-and-ridge}:
\begin{enumerate}
\item \emph{Expressive depth versus compactness:} Raising staircase depth from
two (\(74\) errors) to four in tensor coordinates (\(59\) errors, \(98.00\%\))
essentially matches \emph{compact368} (\(58\) errors, \(98.03\%\)---a single
test window, which is descriptively tied), using \(440\) spectral columns
(\(708\) total) in place of \(100\) (\(368\) total). Continuing to all \(1{,}040\)
unprojected depth-five tensor columns slightly increases errors to \(64\) under
the fixed isotropic penalty \(\alpha = 10\).
\item \emph{Fixed-endpoint ridge sensitivity:} At depth three, the staircase
logsignature (\(60\) spectral columns, \(65\) errors) and staircase tensor
signature (\(160\) spectral columns, \(64\) errors) span the \emph{exact same}
affine linear function space by~\eqref{eq:fixed-endpoint-span}, yet yield
slightly different ridge fits because redundant tensor words alter the singular
values of the standardized feature matrix (Subsection~\ref{app:fixed-endpoints}).
\item \emph{Numerical verification of the staircase projection.} The archived
self-check compares the projected and directly computed QR-Pascal coordinates
on synthetic spectra with absolute tolerance \(10^{-7}\) and relative tolerance
\(10^{-6}\). For the final row of Table~\ref{tab:ordinary-spectral}, the archived
run reports \emph{identical class predictions on all \(2{,}947\) test windows}
after replacing the \(100\) spectral-shape features by the projection. The
mathematical identity holds for all spectra by Theorem~\ref{prop:staircase-projection}.
\end{enumerate}
\FloatBarrier
